\section{Giới thiệu}

Phân đoạn ngữ nghĩa là bài toán gán nhãn cho từng pixel trong ảnh đầu vào, trong đó mô hình phải đồng thời nắm bắt ngữ cảnh toàn cục và bảo toàn các tín hiệu cục bộ. Về mặt hình thức, với ảnh đầu vào
\[
\mathbf{X} \in \mathbb{R}^{H \times W \times 3},
\]
mục tiêu là dự đoán bản đồ nhãn
\[
\mathbf{Y} \in \{1,2,\dots,K\}^{H \times W},
\]
trong đó $H, W$ là kích thước ảnh và $K$ là số lớp cần phân đoạn. Thiết lập này đặc biệt quan trọng trong các hệ thống theo dõi dinh dưỡng, bếp thông minh, robot chế biến và kiểm soát chất lượng thực phẩm, nơi yêu cầu về độ chính xác phải đi kèm với khả năng suy luận hiệu quả trong thực tế \cite{mahalik2010trends, stemmer2018vision, gao2019smartfridge, wang2019smartrefrigerator}.

Trong miền ảnh thực phẩm, bài toán trở nên khó hơn đáng kể vì mỗi pixel không chỉ cần được gán đúng nhãn ngữ nghĩa, mà còn phải được tách khỏi các thành phần nguyên liệu lân cận có hình thái và màu sắc gần nhau. Các thành phần này thường chồng lấp, có ranh giới mờ, biến đổi mạnh theo cách chế biến và mang tính fine-grained rõ rệt. Do đó, chất lượng phân đoạn không chỉ phụ thuộc vào khả năng mô hình hóa ngữ cảnh, mà còn phụ thuộc mạnh vào việc giữ được biên, texture và cấu trúc bề mặt như những tín hiệu phân biệt cục bộ then chốt.

FoodSeg103 là benchmark tiêu biểu cho thiết lập ingredient-level này \cite{benchmarkfood}. Tuy nhiên, đây cũng là một benchmark khó: số lớp lớn, biến thiên nội lớp cao, tương đồng liên lớp mạnh, phân bố đuôi dài nghiêm trọng và mỗi ảnh thường chứa nhiều nguyên liệu cùng lúc. Những đặc tính đó khiến mô hình dễ thiên lệch về lớp phổ biến, làm mịn ranh giới và bỏ sót các lớp hiếm, đặc biệt khi kiến trúc không đủ nhạy với local detail.

Vấn đề nghiên cứu vì thế nằm ở sự đánh đổi giữa độ chính xác và hiệu quả suy luận. Các mô hình segmentation mạnh giàu ngữ cảnh có thể cho kết quả tốt hơn nhưng thường đòi hỏi chi phí tính toán cao, khó triển khai trên thiết bị tài nguyên hạn chế. Ngược lại, các kiến trúc lightweight đạt tốc độ cao nhờ downsample sớm và nén biểu diễn, nhưng chính chiến lược đó có thể làm suy giảm các đặc trưng cục bộ quan trọng đối với ảnh thực phẩm fine-grained. Trong bối cảnh đó, nghiên cứu hiện tại được định vị như một \textbf{baseline study}: thay vì tuyên bố đã giải quyết hoàn chỉnh bài toán texture-aware lightweight segmentation, báo cáo tập trung xây dựng một nền tảng thực nghiệm nhất quán thông qua tái cấu trúc dữ liệu để giảm tác động của long-tail imbalance và chuẩn hóa baseline BiSeNet làm mốc tham chiếu cho các cải tiến texture/local-detail-aware ở giai đoạn tiếp theo.

Từ định vị đó, báo cáo tập trung vào ba đóng góp chính:
\begin{itemize}[leftmargin=1.5em]
    \item Phân tích định lượng bộ dữ liệu FoodSeg103 để chỉ ra các khó khăn cốt lõi của bài toán, bao gồm long-tail imbalance, chồng lấp thành phần, ranh giới yếu và chênh lệch độ phân giải giữa các split.
    \item Xây dựng quy trình làm sạch và tái cấu trúc không gian nhãn dựa trên ngữ nghĩa và tương đồng texture, qua đó giảm số lớp mục tiêu từ 103 xuống 75 trong khi vẫn giữ gần như toàn bộ dữ liệu hợp lệ.
    \item Thiết lập và chuẩn hóa baseline BiSeNet trên bộ dữ liệu đã xử lý để tạo mốc tham chiếu công bằng cho các nghiên cứu texture-aware lightweight segmentation tiếp theo.
\end{itemize}

\subsection{Phát biểu bài toán}

Bài toán phân đoạn ngữ nghĩa (\textit{semantic segmentation}) là bài toán gán nhãn ngữ nghĩa cho từng điểm ảnh trong ảnh đầu vào. Với một ảnh đầu vào
\[
\mathbf{X} \in \mathbb{R}^{H \times W \times 3},
\]
mục tiêu là dự đoán một bản đồ nhãn
\[
\mathbf{Y} \in \{1,2,\dots,K\}^{H \times W},
\]
trong đó $H, W$ lần lượt là chiều cao và chiều rộng của ảnh, còn $K$ là số lớp cần phân đoạn.

Trong bài toán phân đoạn ảnh thực phẩm ở mức nguyên liệu, mỗi pixel cần được gán vào một lớp nguyên liệu tương ứng, ví dụ như \textit{rice}, \textit{beef}, \textit{lettuce}, \textit{tomato}, \dots Với mô hình học sâu, mạng nhận ảnh đầu vào $\mathbf{X}$ và sinh ra tensor logit
\[
\mathbf{Z} \in \mathbb{R}^{H \times W \times K},
\]
trong đó $\mathbf{Z}_{i,j,c}$ biểu diễn điểm số chưa chuẩn hóa của lớp $c$ tại pixel $(i,j)$. Xác suất dự đoán cho lớp $c$ tại pixel $(i,j)$ được tính bằng hàm softmax:
\[
p_{i,j,c} = \frac{\exp(\mathbf{Z}_{i,j,c})}{\sum_{k=1}^{K} \exp(\mathbf{Z}_{i,j,k})}.
\]

Nhãn dự đoán cuối cùng tại pixel $(i,j)$ được xác định bởi:
\[
\hat{y}_{i,j} = \arg\max_{c \in \{1,\dots,K\}} p_{i,j,c}.
\]

Như vậy, bài toán phân đoạn ngữ nghĩa có thể xem là bài toán phân lớp đa lớp ở mức pixel, trong đó mô hình phải đồng thời học được ngữ cảnh toàn cục và bảo toàn các chi tiết cục bộ như biên, texture và cấu trúc bề mặt để phân biệt các lớp có mức độ tương đồng thị giác cao.

\subsection{Hàm mất mát}

Trong nghiên cứu này, mô hình được huấn luyện bằng hàm mất mát \textit{pixel-wise cross-entropy}. Gọi $y_{i,j}$ là nhãn thật tại pixel $(i,j)$, khi đó hàm mất mát trên toàn ảnh được viết là:
\[
\mathcal{L}_{\mathrm{CE}}
=
-\frac{1}{|\Omega|}
\sum_{(i,j)\in\Omega}
\log p_{i,j,y_{i,j}},
\]
trong đó $\Omega$ là tập các pixel hợp lệ được dùng để huấn luyện, và $|\Omega|$ là số lượng pixel trong tập này.

Viết tường minh theo đầu ra softmax, ta có:
\[
\mathcal{L}_{\mathrm{CE}}
=
-\frac{1}{|\Omega|}
\sum_{(i,j)\in\Omega}
\log
\left(
\frac{\exp(\mathbf{Z}_{i,j,y_{i,j}})}
{\sum_{c=1}^{K}\exp(\mathbf{Z}_{i,j,c})}
\right).
\]

Hàm mất mát này buộc mô hình tăng xác suất cho lớp đúng tại từng pixel. Trong thực tế, nếu dữ liệu có mất cân bằng lớp mạnh, có thể mở rộng sang \textit{weighted cross-entropy}; tuy nhiên, dạng cơ sở được sử dụng trong báo cáo này là cross-entropy chuẩn ở mức pixel.

\subsection{Các chỉ số đánh giá}

Để đánh giá chất lượng phân đoạn, ta xét theo từng lớp $c$. Gọi:
\begin{itemize}[leftmargin=1.5em]
    \item $TP_c$ là số pixel thuộc lớp $c$ và được dự đoán đúng là lớp $c$;
    \item $FP_c$ là số pixel không thuộc lớp $c$ nhưng bị dự đoán nhầm thành lớp $c$;
    \item $FN_c$ là số pixel thuộc lớp $c$ nhưng bị dự đoán nhầm sang lớp khác.
\end{itemize}

Khi đó, \textit{Intersection over Union} của lớp $c$ được định nghĩa là:
\[
\mathrm{IoU}_c = \frac{TP_c}{TP_c + FP_c + FN_c}.
\]

\subsection{Mean IoU trên các lớp xuất hiện}

Chỉ số đánh giá chính trong nghiên cứu này là \textbf{mean Intersection over Union (mIoU)}. Tuy nhiên, thay vì lấy trung bình trên toàn bộ $K$ lớp cố định, mIoU được tính trên \textbf{các lớp thực sự xuất hiện trong tập đánh giá}. Đây là điểm cần nhấn mạnh để tránh làm sai lệch kết quả trong bối cảnh dữ liệu có phân bố lớp thưa và không phải lớp nào cũng xuất hiện trong mỗi pha đánh giá.

Gọi tập các lớp xuất hiện là:
\[
\mathcal{C}_{\mathrm{present}}
=
\left\{
c \in \{1,\dots,K\}
\;\middle|\;
TP_c + FN_c > 0
\right\}.
\]
Tức là một lớp được xem là \textit{xuất hiện} nếu nó có ít nhất một pixel ground-truth trong tập đánh giá.

Khi đó, mIoU được tính như sau:
\[
\mathrm{mIoU}
=
\frac{1}{|\mathcal{C}_{\mathrm{present}}|}
\sum_{c \in \mathcal{C}_{\mathrm{present}}}
\mathrm{IoU}_c.
\]

Cách tính này phản ánh đúng hơn chất lượng phân đoạn trên các lớp thực sự có mặt, thay vì đưa cả những lớp không xuất hiện vào mẫu số trung bình. Do đó, đây là chỉ số phù hợp hơn cho các bài toán phân đoạn thực phẩm fine-grained, nơi phân bố lớp thường mất cân bằng và không đồng đều giữa các mẫu.

\subsection{Chỉ số bổ sung}

Ngoài mIoU, có thể báo cáo thêm một số chỉ số bổ sung để hỗ trợ phân tích:

\paragraph{Mean Accuracy (mAcc).}
Độ chính xác theo từng lớp được định nghĩa là:
\[
\mathrm{Acc}_c = \frac{TP_c}{TP_c + FN_c}.
\]
Tương tự, trung bình trên các lớp xuất hiện là:
\[
\mathrm{mAcc}
=
\frac{1}{|\mathcal{C}_{\mathrm{present}}|}
\sum_{c \in \mathcal{C}_{\mathrm{present}}}
\mathrm{Acc}_c.
\]

\paragraph{Overall Accuracy (aAcc).}
Độ chính xác toàn cục trên tất cả pixel:
\[
\mathrm{aAcc}
=
\frac{\sum_{c} TP_c}{\sum_{c}(TP_c + FN_c)}.
\]

Tuy nhiên, trong báo cáo này, \textbf{mIoU được sử dụng làm chỉ số chính} vì đây là thước đo phù hợp hơn cho bài toán phân đoạn ngữ nghĩa nhiều lớp, đặc biệt trong trường hợp dữ liệu mất cân bằng giữa các lớp.